\chapter{Control, Error Budgets and Accuracy Standards}

\section{Relative versus absolute}

\begin{kid}
You can draw a perfect map of your bedroom --- every toy in the right place compared to every other toy --- but if you don't know where your house is on the street, nobody else can use it to find your room. SLAM makes the bedroom map. \textbf{Control points} --- spots measured very carefully by surveyors --- pin that map onto the real world. \textbf{Checkpoints} are extra careful spots that we keep secret from the computer, then use to test whether the map is truly right.
\end{kid}

\fig[0.82]{control}{Along a corridor, alternate points used to constrain the solution (control) with points held back purely to test it (checkpoints).}

\begin{tip}
\begin{itemize}
\item Use targets the scanner can see and the software can measure: painted chevrons or checkerboards on pavement for vertical and horizontal checks; high-intensity paint shows clearly in intensity.
\item Survey control and checkpoints with a method at least \textbf{two times} more accurate than the product you are testing (ASPRS Edition 2), and include their own uncertainty in the final accuracy (below).
\item Distribute checkpoints along the whole project and across the cross-section --- not only at the ends near GNSS-friendly areas.
\item Confirm the vertical datum and geoid model match project requirements before comparing anything.
\end{itemize}
\end{tip}

\section{Computing accuracy from checkpoints}

\fig[0.78]{residuals}{Vertical residuals (LiDAR surface minus surveyed elevation) at ten checkpoints.}

\section{ASPRS lidar data standards in detail}

\begin{kid}
One rulebook says \emph{how close to the truth} the dots must be and how to check them with secret checkpoints. The other rulebook says \emph{how to write the dots down} in a file, including a label on every dot: ground, building, tree, water, wire, noise\ldots{} so any software can read it.
\end{kid}
